\scanpage{89}
% L1–5
La formule \textcircled{\scriptsize 10} est \unsure{importante} pour le calcul numérique des fonctions $\Theta$, car ramène le calcul de $\Theta(v\mid iT)$ à celui de $\Theta(v\mid\frac iT)$; si donc $\unsure{\Re|\mathcal R}T\ngtr1$, alors on a $\Re T>1$ ou $\Re\frac1T>1$, d'où résulte que dans l'expression

% L6
\[
\Theta(v\mid iT)=2\sum_1^\infty(-1)^{n-1}h^{(n-\frac1{\unsure{2|3}})^2}\sin(2n-1)\pi v,\quad\text{où}
\]

% L7–11
$h=e^{-\pi T}$ ou $e^{-\pi\frac1T}$, on peut supposer que $|h|\leq e^{-\pi}\mathrel{\#}\frac1{10}$; comme l'exposant de $h$ augmente rapidement, on voit qu'on a bien une série très rapidement convergente.

% L12, underlined
\uline{Autres fonctions $\Theta$ liées aux fonctions elliptiques}

% L13
La fonction $\Theta_i$ $(i=1,2,3)$ est définie par

% L14
\[
\sigma_i u=e^{\frac{\eta_1u^2}{2\omega_1}}\,
\frac{\Theta_i\left(u/2\omega_1\,\middle|\,\frac{\omega_2}{\omega_1}\right)}
{\Theta_i\left(0\,\middle|\,\frac{\omega_2}{\omega_1}\right)}
\tag*{\textcircled{\scriptsize 13}}
\]

% L15
et s'exprime par une série du type $\Theta$ :

% L16–19, brace continuing
\[
\left\{\begin{aligned}
\Theta_1(v\mid t)&=\deleted{\unsure{2|e}}\sum_{-\infty}^\infty e^{(n-\frac1{\unsure{1|2}})^2i\pi t+(2n-1)i\pi v}\\
&=2\sum_1^\infty e^{(n-\frac11)^2i\pi t}\cos(2n-1)\pi v=\vartheta_2(v)\\
\Theta_2(v\mid t)&=\sum_{-\infty}^{+\infty}(\unsure{-1|-i})^n e^{n^2i\pi t+2ni\pi v}\deleted{\unsure{\sum^{\infty}|2\sum}}\\
&=1+2\sum_1^\infty e^{n^2i\pi t}\cos2\unsure{\pi|n}\pi v=\vartheta_0(v)
\end{aligned}\right.
\tag*{\textcircled{\scriptsize 14}}
\]

% lower left
% archive sheet number: 88
